% !PRINTED 232
\weberpage{232}
%<<CONT
worin $\left(\frac{2}{n}\right) = (-1)^{\frac{n^2-1}{8}}$, und wenn wir also die dem Falle $\mu = 1$, $\mu' = 0$ entsprechenden Funktionen $P$ mit $P^0$ bezeichnen, so folgt aus § 67, (2) und (4) mit Rücksicht auf § 34, (10):
\weq{1}{\begin{aligned}
P_{00}^0 &= \frac{f(n\omega)}{f(\omega)^n},\\
P_{01}^0 &= \frac{f_1(n\omega)}{f_1(\omega)^n},\\
P_{10}^0 &= \left(\frac{2}{n}\right)\frac{f_2(n\omega)}{f_2(\omega)^n},
\end{aligned}}
\weq{2}{P_{11}^0 = \sqrt{n}\,\frac{\eta(n\omega)}{\eta(\omega)}.}

Um durch solche Formeln die sämtlichen Wurzeln der Transformationsgleichungen darzustellen, bestimmen wir eine lineare Transformation folgendermaßen:
\weq{3}{\begin{pmatrix}\alpha, & \beta\\ \gamma, & \delta\end{pmatrix} \equiv \begin{pmatrix}1, & 0\\ 0, & 1\end{pmatrix} \pmod{16},}
\weq{4}{\alpha \equiv \mu, \qquad \beta \equiv \mu' \pmod{n},}
und ersetzen $\omega$ in (1), (2) durch
\weq{5}{\omega' = \frac{\gamma + \delta\omega}{\alpha + \beta\omega};}
auf die linke Seite lassen sich dann die Formeln (3) bis (6) § 39 und (19), § 38 anwenden und ergeben:
\weq{6}{\begin{aligned}
P_{00} &= \frac{f(n\omega')}{f(\omega')^n},\\
P_{01} &= \frac{f_1(n\omega')}{f_1(\omega')^n},\\
P_{10} &= \left(\frac{2}{n}\right)\frac{f_2(n\omega')}{f_2(\omega')^n},\\
P_{11} &= \varepsilon^{1-n}\sqrt{n}\,\frac{\eta(n\omega')}{\eta(\omega')},
\end{aligned}}
worin $\varepsilon$ die in § 33, (15), (18) genau definierte 24ste Einheitswurzel ist.

Nun lassen sich zwei Transformationen
\[
% [sic] The source prints a semicolon after \alpha' in the first matrix.
\begin{pmatrix}\alpha'; & \beta'\\ \gamma', & \delta'\end{pmatrix}, \quad \begin{pmatrix}a, & 0\\ c, & \partial\end{pmatrix}, \quad a\partial = n,
\]
%CONT>>
